\documentclass[11pt]{article}
\newcommand{\paso}{7}
\usepackage[T1]{fontenc}
\usepackage[utf8]{inputenc}
\usepackage[spanish,es-nodecimaldot,es-noquoting,es-lcroman]{babel}
\spanishplainpercent
\usepackage{lmodern}
\usepackage[a4paper,margin=21mm,headheight=15pt]{geometry}
\usepackage{amsmath,amssymb,bm}
\usepackage{booktabs,array}
\usepackage{tikz}
\usetikzlibrary{arrows.meta,positioning,calc}
\usepackage[most]{tcolorbox}
\usepackage{microtype}
\usepackage{fancyhdr}
\usepackage{hyperref}
\hypersetup{colorlinks=true,linkcolor=blue!45!black,urlcolor=blue!45!black}
\setlength{\parindent}{0pt}
\setlength{\parskip}{5pt}
\setlength{\emergencystretch}{2em}
\pagestyle{fancy}
\fancyhf{}
\fancyhead[L]{\small BEM del cilindro: paso \paso}
\fancyhead[R]{\small Guía matemática independiente}
\fancyfoot[C]{\thepage}
\newcommand{\vct}[1]{\bm{#1}}
\newcommand{\dd}{\,\mathrm{d}}
\newcommand{\ii}{\mathrm{i}}
\newcommand{\norm}[1]{\left\lVert#1\right\rVert}
\newcommand{\VP}{\operatorname{VP}}
\newcommand{\tr}{\operatorname{tr}}
\newtcolorbox{idea}[1]{breakable,colback=blue!3,colframe=blue!50!black,
  title={#1},fonttitle=\bfseries,boxrule=0.7pt}
\newtcolorbox{alerta}[1]{breakable,colback=orange!4,colframe=orange!65!black,
  title={#1},fonttitle=\bfseries,boxrule=0.7pt}
\newtcolorbox{resultado}[1]{breakable,colback=teal!3,colframe=teal!55!black,
  title={#1},fonttitle=\bfseries,boxrule=0.7pt}
\newcommand{\inicio}[2]{%
 {\small\color{blue!55!black}\bfseries SERIE: LOS SIETE PASOS MATEMÁTICOS DEL BEM}\par
 {\LARGE\bfseries Paso \paso. #1\par}\vspace{3mm}
 \begin{idea}{¿Qué busca este paso?}#2\end{idea}}
\newcommand{\mapa}[3]{%
 \begin{center}
 \begin{tikzpicture}[node distance=4mm,
   caja/.style={draw=blue!50!black,fill=blue!3,rounded corners=2pt,
     text width=13cm,align=center,inner sep=6pt},
   >={Stealth[length=2mm]}]
 \node[caja] (a){#1};
 \node[caja,below=of a] (b){#2};
 \node[caja,below=of b] (c){#3};
 \draw[->,thick,blue!60!black] (a)--(b);
 \draw[->,thick,blue!60!black] (b)--(c);
 \end{tikzpicture}
 \end{center}}
\newcommand{\base}{%
 \smallskip{\footnotesize\color{black!65}
 Modelo de referencia: cilindro sólido, material homogéneo e isótropo,
 deformación plana, pequeñas deformaciones, sin fuerzas de volumen;
 amplitudes complejas con \(e^{+\ii\omega t}\), \(\omega>0\).}\par}
\newcommand{\fuentes}{%
 \par\medskip\begingroup\footnotesize
 \textbf{Referencia y continuación.}
 Este documento desarrolla una etapa del esquema
 \texttt{esquema\_ecuaciones\_bem.tex}. Las ecuaciones y los ejemplos se
 contrastaron con \texttt{bem\_cilindro.py}, ambos en la carpeta superior.
 La guía general y sus deducciones ampliadas se encuentran en
 \texttt{../guia\_bem/guia\_bem.pdf}. No se modifica la implementación.
 \par Para recompilar este documento, mantener \texttt{estilo.tex} en
 la misma carpeta y ejecutar \texttt{pdflatex} sobre su archivo
 \texttt{.tex}, desde esa carpeta.\par\endgroup}

\begin{document}
\inicio{Aplicar, resolver e interpretar}{Usar la presión conocida para
construir el segundo miembro, encontrar los desplazamientos nodales
complejos y convertirlos en una respuesta radial de amplitud y fase.
El último objetivo es distinguir una solución algebraica correcta de
una predicción física numéricamente fiable.}
\base
\mapa{Datos: \(H(\omega),G(\omega)\), presión \(p\), nodos \(\vct x_g\)}
{Carga nodal: \(\vct t_g=-p\,\vct x_g/R\);\quad resolver \(H\vct u=G\vct t\)}
{Salida: \(u_{r,g}=\vct u_g\cdot\vct n_g\), amplitud \(|u_{r,g}|\) y fase \(\arg u_{r,g}\)}

\section*{1. Construir el vector conocido}
En cada nodo circular,
\[
 \vct n_g=\frac{\vct x_g}{R}=(\cos\theta_g,\sin\theta_g),\qquad
 \vct t_g=-p(\cos\theta_g,\sin\theta_g).
\]
Después se apilan las componentes en el mismo orden del ensamblaje:
\[
 \boxed{\vct t=
 (t_{0x},t_{0y},t_{1x},t_{1y},\ldots,t_{N-1,x},t_{N-1,y})^{\mathsf T}.}
\]
Este orden debe coincidir con las columnas de \(G\). Para ocho elementos
hay \(N=16\) y \(\vct t\) tiene 32 entradas.

\section*{2. El segundo miembro no es la presión}
\[
 \boxed{\vct b=G(\omega)\vct t,\qquad H(\omega)\vct u=\vct b.}
\]
\(\vct t\) está en Pa; \(\vct b\), en metros. Se interpreta como el término
conocido de influencia de las tracciones, no como una fuerza nodal FEM.
La integración sobre el contorno ya está contenida en \(G\).

\begin{alerta}{El signo de la compresión}
En \((R,0)\), \(\vct n=(1,0)\), por lo que
\(\vct t=(-p,0)\). Usar \(+p\vct n\) representa tracción hacia afuera.
Una vez elegido el signo de la presión, no se introduce otro signo menos
durante la solución.
\end{alerta}

\newpage
\section*{3. Resolver el sistema complejo}
El sistema contiene las dos componentes cartesianas del desplazamiento de
cada nodo:
\[
 \vct u=(u_{0x},u_{0y},u_{1x},u_{1y},\ldots)^{\mathsf T}.
\]
Formalmente se podría escribir \(H^{-1}\vct b\), pero no conviene construir
la inversa. Se emplea una resolución de sistema, por ejemplo factorización
LU con pivoteo:
\[
 PH=LU,\qquad L\vct y=P\vct b,\qquad U\vct u=\vct y.
\]
Aquí \(P\) es una matriz de permutación, no la presión \(p\).
La matriz triangular \(U\) de esta factorización tampoco es el núcleo
fundamental; la letra se usa en un contexto algebraico diferente.

\section*{4. Qué puede salir mal al resolver}
\begin{itemize}
 \item Una matriz singular no produce una solución única mediante una
 resolución directa ordinaria.
 \item Una matriz mal condicionada amplifica errores en el operador o la
 carga. Cerca de una resonancia, la respuesta física puede ser grande.
 \item A \(\omega=0\), un problema estático de tracciones tiene movimientos
 rígidos indeterminados. La implementación dinámica se usa con
 \(\omega>0\), no sustituyendo cero en las fórmulas de Bessel.
 \item Para resonancias o problemas espectrales delicados se necesita
 analizar el operador y su discretización; añadir amortiguamiento no es
 una prueba universal de exactitud.
\end{itemize}

\section*{5. Dos cargas y dos frecuencias no se tratan igual}
Si geometría, material y frecuencia no cambian, \(H,G\) son los mismos para
distintas presiones. Por linealidad:
\[
 \vct u(\alpha p)=\alpha\,\vct u(p).
\]
Se puede reutilizar la factorización para varias cargas a la misma
frecuencia. Si cambia \(\omega\), hay que reconstruir \(H,G\) y factorizar
de nuevo; solo el precálculo geométrico y estático permanece reutilizable.

\begin{resultado}{Control algebraico}
\[
 r_{\mathrm{alg}}=
 \frac{\norm{H\vct u-G\vct t}}{\norm{G\vct t}}.
\]
Esta razón es útil si el denominador es distinto de cero. Para una carga
nula se usa un control absoluto o una normalización adecuada.
Un residuo pequeño certifica la resolución del sistema ensamblado,
no la exactitud de la discretización física.
\end{resultado}

\newpage
\section*{6. Proyectar a la dirección radial}
En el nodo \(g\), las bases radial y tangencial son
\[
 \vct e_{r,g}=(\cos\theta_g,\sin\theta_g)=\vct n_g,\qquad
 \vct e_{\theta,g}=(-\sin\theta_g,\cos\theta_g).
\]
Los desplazamientos son
\[
 \boxed{u_{r,g}=u_{x,g}\cos\theta_g+u_{y,g}\sin\theta_g,}
\]
\[
 u_{\theta,g}=-u_{x,g}\sin\theta_g+u_{y,g}\cos\theta_g.
\]
La proyección usa las componentes complejas, sin tomar valores absolutos
ni conjugados previamente. La normal geométrica del nodo circular es real.
Por simetría, el problema exacto tiene \(u_\theta=0\) y el mismo \(u_r\)
en todos los nodos superficiales.

\section*{7. De una amplitud compleja al movimiento real}
Si \(u_r=a+\ii b=Ae^{\ii\phi}\), entonces
\[
 A=\sqrt{a^2+b^2},\qquad \phi=\operatorname{atan2}(b,a),
\]
\[
 \boxed{u_{r,\text{fís}}(t)=a\cos\omega t-b\sin\omega t
 =A\cos(\omega t+\phi).}
\]
\(A\ge0\) es amplitud, no contracción con signo. El signo y el desfase
están en \(\phi\) o en las partes real e imaginaria.

\begin{resultado}{Ejemplo a \(\omega=40\) rad/s, ocho elementos}
Con el material y presión de referencia, en el nodo 0:
\[
 u_{r,0}\simeq-2.395966\times10^{-4}
                +\ii\,4.085472\times10^{-5}\ \mathrm m.
\]
\[
 |u_{r,0}|\simeq2.430548\times10^{-4}\ \mathrm m
 =0.243055\ \mathrm{mm},\qquad
 \phi_0\simeq170.3233^\circ.
\]
En \(t=0\), la posición radial relativa es la parte real negativa,
no la amplitud positiva. La fase indicada está referida a la presión
escalar positiva \(p\cos\omega t\); la tracción radial prescrita es \(-p\).
\end{resultado}

\newpage
\section*{8. Comprobar con una solución analítica independiente}
En la respuesta axisimétrica exacta, sea
\(\vct u=u_r(r)\vct e_r\). La ecuación radial se reduce a
\[
 u_r''+\frac1r u_r'-\frac{u_r}{r^2}+k_p^2u_r=0,\qquad
 k_p=\omega/c_p.
\]
La solución regular en el centro es
\[
 u_r(r)=A J_1(k_pr).
\]
Se descarta \(Y_1\) por su singularidad en el origen.
Aquí \(J_0,J_1\) son funciones de Bessel de primera especie:
no son el jacobiano geométrico ni las \(K_n\) de los núcleos auxiliares.

\subsection*{8.1 Imponer la presión en \(r=R\)}
La ley constitutiva en deformación plana da
\[
 \sigma_{rr}=(\lambda^*+2\mu^*)u_r'
                  +\lambda^*\frac{u_r}{r}.
\]
Usando \(J_1'(z)=J_0(z)-J_1(z)/z\):
\[
 \sigma_{rr}(R)=A\left[
 (\lambda^*+2\mu^*)k_pJ_0(k_pR)-\frac{2\mu^*}{R}J_1(k_pR)
 \right]=-p.
\]
Si \(x=k_pR\), resulta
\[
 \boxed{u_r^{\mathrm{exacto}}(R)=
 -\frac{pR\,J_1(x)}
 {(\lambda^*+2\mu^*)xJ_0(x)-2\mu^*J_1(x)}.}
\]
La comparación comprueba la geometría, la carga, las ondas y la resolución,
no solo el álgebra del sistema.

\subsection*{8.2 Límite de baja frecuencia}
Con \(J_1(x)\simeq x/2\) y \(J_0(x)\simeq1\):
\[
 u_r(R)\longrightarrow-\frac{pR}{2(\lambda^*+\mu^*)}.
\]
Para el material amortiguado de referencia:
\[
 u_r(R)\longrightarrow
 -1.485149\times10^{-4}+\ii1.485149\times10^{-5}\ \mathrm m.
\]
Es un límite de la solución física. No implica que sea seguro evaluar
los núcleos dinámicos directamente con \(\omega=0\).

\newpage
\section*{9. Promedios: qué magnitud se está comparando}
La clase devuelve los \(N\) valores complejos \(u_{r,g}\), no una media.
Dos resúmenes diferentes son
\[
 \overline A=\frac1N\sum_g|u_{r,g}|,\qquad
 A_{\mathrm{coh}}=\left|\frac1N\sum_gu_{r,g}\right|.
\]
En general, \(A_{\mathrm{coh}}\le\overline A\). Por ejemplo, dos respuestas
\(a\) y \(-a\) tienen amplitud media \(|a|\), pero promedio complejo cero.
En una respuesta perfectamente axisimétrica ambas medidas coinciden.
Estos son promedios nodales; no son automáticamente integrales ponderadas
por longitud de arco.

\section*{10. Resultado numérico y qué prueba cada control}
Para \(R=6\) m, \(p=100\) Pa, \(\beta=0.05\), ocho elementos,
\(\omega=40\) rad/s:
\[
 \overline A_{\mathrm{BEM}}\simeq2.430943\times10^{-4}\ \mathrm m,\qquad
 A_{\mathrm{exacta}}\simeq2.433336\times10^{-4}\ \mathrm m.
\]
\[
 e_A=\frac{|\overline A_{\mathrm{BEM}}-A_{\mathrm{exacta}}|}
                {A_{\mathrm{exacta}}}
 \simeq9.8369\times10^{-4}=0.09837\%.
\]
\begin{itemize}
 \item \textbf{Álgebra:} el residuo relativo observado es aproximadamente
 \(2.5\times10^{-16}\). No equivale al error físico anterior.
 \item \textbf{Simetría:} revisar uniformidad angular de \(u_r\) y que
 \(u_\theta\) sea pequeño; se observó
 \(\max_g|u_{\theta,g}|\simeq5.0\times10^{-10}\) m.
 \item \textbf{Fase y campo complejo:} un error pequeño de amplitud no
 garantiza fase correcta. Comparar también \(u_r\) complejo.
 \item \textbf{Convergencia:} repetir con más elementos y controlar la
 cuadratura por separado. Un buen resultado en una frecuencia no certifica
 todo el barrido.
\end{itemize}
\begin{alerta}{No confundir tres errores}
El error de resolución lineal, el error de integración y el error de
representación geométrica y de campos son distintos. Pueden coexistir un
residuo diminuto y una respuesta física apreciablemente equivocada.
\end{alerta}
\section*{11. Cierre del flujo matemático}
El producto final es un desplazamiento radial complejo por nodo y frecuencia.
Al variar \(\omega\) se obtiene una curva de respuesta; al variar \(p\) en
el régimen lineal se escala la solución.
\fuentes
\end{document}
